\documentclass[10pt,a4paper]{amsart}
\usepackage[margin=19mm]{geometry}
\usepackage{amsmath,amssymb,mathtools}
\usepackage[scr=rsfs,cal=euler]{mathalfa}
\usepackage{xfrac}
\usepackage[unicode,hidelinks,bookmarks=false]{hyperref}
\newcommand{\ie}{i.e.\ }
\newcommand{\cf}{cf.\ }
\newcommand{\resp}{respectively\ }
\newcommand{\Subsection}{\subsection}
\providecommand{\nobreakdash}{\nobreak\hbox{-}}
\setlength{\emergencystretch}{2em}
\multlinegap0pt
\usepackage{bm}
\usepackage{bbm}
\usepackage{dsfont}
\usepackage{xspace}
\usepackage{cancel}
\usepackage{mathtools}
\usepackage{amsthm}
\makeatletter
\@ifundefined{definition}{\newtheorem{definition}{Definition}}{}
\@ifundefined{lemma}{\newtheorem{lemma}{Lemma}}{}
\makeatother
\DeclareMathOperator{\distG}{\operatorname{dist_\mathcal{G}}} % graph distance 
\def\N{\mathbb{N}} 
\newcommand{\G}{\mathcal{G}}
\newcommand{\Rn}{\R^n}
\newcommand{\R}{\mathbb{R}}
\newcommand{\abs}[1]{\left\lvert {#1} \right\rvert}
\newcommand{\agents}{\lbrace 1, \dots, s \rbrace}
\newcommand{\norm}[1]{\left\lVert {#1} \right\rVert}
\renewcommand{\L}{\mathcal{L}}
\title[Separable Approximations of Optimal Value Functions and Thei]{Separable Approximations of Optimal Value Functions and Their Representation by Neural Networks}
\author{Mario Sperl, Luca Saluzzi, Dante Kalise, Lars Gr\"une}
\date{Source: arXiv:2502.08559v2 (CC BY 4.0)}
\begin{document}
\maketitle

\noindent\emph{Source and licence.} Separable Approximations of Optimal Value Functions and Their Representation by Neural Networks. Version arXiv:2502.08559v2 (CC BY 4.0), distributed under the Creative Commons Attribution 4.0 International licence, as recorded in the arXiv licence metadata for this exact version and shown on the abstract page https://arxiv.org/abs/2502.08559v2. Full rights notice: licenses/arxiv-2502.08559-CC-BY-4.0.txt.\par\smallskip

\noindent\textit{Editorial scope.} Counted unit: the Lemma whose source label is \texttt{lem:SepApprox:QuadraticEstimate} in the article (source0.tex lines 295--307, source label \texttt{lem:SepApprox:QuadraticEstimate}), delivered together with the article's own complete original proof of it (source0.tex lines 308--346, one \texttt{proof} environment of 39 lines). No mathematics is written, repaired or re-derived: statement and proof are the article's own text line for line. Required context retained uncounted: eq:setting:def\_gj (lines 166--169); def:setting:sensitivity (lines 174--179); eq:SepApprox:DefPsij (lines 252--254). Every definition, display and label of those lines is retained unchanged. Omitted from this package and not redistributed: 1--165, 170--173, 180--251, 255--294, 347--1140. Nothing inside a retained excerpt is omitted or altered: every retained body is delivered exactly as the article's lines are written, so no omission receipt of any kind accompanies this package. Citations: the retained excerpts carry no citation command at all, so no bibliography entry of the article is transcribed and no reference is redistributed. Reference adaptation: none was needed and none was made; every reference inside the delivered unit resolves inside it, so no unresolved reference or citation is produced. Printed slips kept literal: no printed word, symbol, hypothesis or number of the counted statement or of its proof was altered. Layout and class: the article's own class is \texttt{article} with packages including geometry, lipsum, amsfonts, graphicx, epstopdf, algorithmic, amsmath, amsthm, amscd, amssymb, mathtools, mathrsfs, babel, inputenc; the wrapper is the lane's pinned \texttt{amsart} editorial wrapper at 10pt on A4, which restates the theorem-like environments the retained text needs and adds the article's own macro definitions that the retained text needs (\texttt{\textbackslash G}, \texttt{\textbackslash L}, \texttt{\textbackslash N}, \texttt{\textbackslash R}, \texttt{\textbackslash Rn}, \texttt{\textbackslash abs}, \texttt{\textbackslash agents}, \texttt{\textbackslash distG}, \texttt{\textbackslash norm}), with extra wrapper packages (bm, bbm, dsfont, xspace, cancel, mathtools). The delivered numbering is the unit's own. Assets: no retained span contains a figure environment, a graphics-inclusion command, a PDF-page inclusion, a table or a TikZ picture; no image file is copied into this package, the package declares no asset dependency and no image file is redistributed with it. Licence: version arXiv:2502.08559v2 of this article is distributed under the Creative Commons Attribution 4.0 International licence, as recorded in the arXiv licence metadata for this exact version and shown on the abstract page https://arxiv.org/abs/2502.08559v2. A full rights notice is delivered at licenses/arxiv-2502.08559-CC-BY-4.0.txt. Characters: every retained excerpt is pure ASCII, so the delivered engine, XeLaTeX with its default Latin Modern text font, reports no missing character.
    \begin{equation} \label{eq:setting:def_gj}
        g_j \colon \Rn \to \R , \qquad 
        x = (z_1, \dots, z_s) \mapsto V(x) - V(z_1, \dots, z_{j-1}, 0, z_{j+1}, \dots, z_s).
    \end{equation}

\begin{definition} \label{def:setting:sensitivity}
   We say that $V$ has the $\gamma$-decaying sensitivity property with respect to $\G$ on $\Omega$ for $\gamma \in \L$ if for all $i, j \in \agents$ with $i \neq j$ and for all $x \in \Omega$ the mapping $g_j$ as defined in \eqref{eq:setting:def_gj} is Lipschitz w.r.t. $z_i$ at the point $x = (z_1, \dots, z_s)$ with Lipschitz constant
    \begin{equation} \label{eq:def_Lipschitz}
        L_{i,j} \coloneqq \gamma (\distG(i,j)) \norm{z_j}. 
    \end{equation}
\end{definition}

\begin{equation} \label{eq:SepApprox:DefPsij}
    \Psi_l^j(z_{j,l}) \coloneqq V(x_{j,l}) - V(\Lambda_j x_{j,l}), 
\end{equation}

\begin{lemma} \label{lem:SepApprox:QuadraticEstimate}
     Assume that $V$ has the $\gamma$-decaying sensitivity property. Then, for any $l \in \N$ and $x \in \Omega$, we have 
    \begin{equation*}
        \lvert V(x) - \Psi_l(x) - V(0) \rvert \leq \lVert x \rVert_2^2 \lVert D_l \rVert_2, 
    \end{equation*}
    where $D_l \in \R^{s \times s}$ is given by 
    \begin{equation}\label{def.D}
        D_l[i,j] := \begin{cases}
        \gamma({\distG(i,j)}), \quad & \text{if} \, \distG(i,j) > l, \\
       0, \quad & \text{otherwise}. 
    \end{cases}
    \end{equation}
\end{lemma}

\begin{proof}
Initially, we observe that for any $i, j \in \agents$ we have that either $\distG(i,j) \leq s-1 $ or $\distG(i,j) = \infty$. In the latter case it follows that $V(x) - V(\Lambda_j x)$ is independent of $z_i$. Thus, by considering graph-neighborhoods of the size $s-1$ we can express $V(x)$ as 
$V(x) = V(0) + \sum_{j=1}^s \Psi_{s-1}^j(z_{j, s-1})$. This shows the claim for $l \ge s-1$. For the case $l \le s-2$ we use a telescoping sum to write  
\begin{align}       \begin{split}\label{eq:SepApprox:TelescopeV}
        & \Big| {V(x) - \sum_{j=1}^s \Psi_l^j(z_{j,l}) - V(0)}  \Big| =  \Big| {\sum_{j=1}^s \Psi_{s-1}^j(z_{j,s-1}) - \sum_{j=1}^s \Psi_l^j(z_{j,l})}  \Big| \\ \leq & \sum_{j=1}^s \sum_{k = l}^{s-2} \abs{\Psi_{k+1}^j(z_{j,k+1}) - \Psi_{k}^j(z_{j,k})}. 
    \end{split}
\end{align}
Now, fix some $j \in \agents$ and $\tilde{l} \in \{l, \dots, s-2\}$. Let $\tau \in \N$ and $i_1, \dots, i_\tau \in \agents$ such that $\lbrace i \in \agents \mid \distG(j,i) = \tilde l \rbrace = \lbrace i_1, \dots, i_\tau \rbrace $.
Then, by the definitions of $\Psi_l^j$ in \eqref{eq:SepApprox:DefPsij} and $g_j$ in \eqref{eq:setting:def_gj} we have
    \begin{align*}
        & \Psi_{\tilde l}^j(z_{j,{\tilde l}}) -\Psi_{{\tilde l}+1}^j(z_{j,{\tilde l}+1}) = V(x_{j,{\tilde l}}) - V(\Lambda_j x_{j,{\tilde l}}) - V(x_{j,{\tilde l}+1}) + V(\Lambda_j x_{j,{\tilde l}+1}) \\ 
        = &  \sum_{p = 1}^\tau g_j(x_{j, {\tilde l}+1}^{(p)}) - g_j(x_{j, {\tilde l}}^{(p)}), 
    \end{align*}
    where for each $1 \leq p \leq \tau$ the vectors $x_{j, {\tilde l}+1}^{(p)}$ and $x_{j, {\tilde l}}^{(p)}$ are obtained by setting all entries corresponding to $z_{i_1}, \dots, z_{i_p}$ to $0$ in the vectors $x_{j, {\tilde l}+1}$ and $x_{j,{\tilde l}}$, respectively. By successively applying the $\gamma$-decaying sensitivity property as defined in Definition \ref{def:setting:sensitivity}, we obtain
    \begin{equation*}
        \abs{ \Psi_{\tilde l+1}^j(z_{j,\tilde l+1}) -\Psi_{\tilde l}^j(z_{j,\tilde l})} \leq \gamma(\tilde l+1) \sum_{p=1}^\tau \norm{z_{i_p}} \norm{z_j}. 
    \end{equation*}
    From this, using \eqref{eq:SepApprox:TelescopeV} we can derive
    \begin{align*}
         & \Bigg| {V(x) - \sum_{j=1}^s \Psi_l^j(z_{j,l}) - V(0)} \Bigg|  \leq \sum_{j=1}^s \sum_{k = l}^{s-2} \; \sum_{i: \distG(i,j) = k+1} \norm{z_i} \gamma(k+1) \norm{z_j} \\ 
         = &  \sum_{j=1}^s \sum_{i: \distG(i,j) > l} \norm{z_i} \gamma(\distG(i,j)) \norm{z_j}
         =  \sum_{j=1}^s \norm{z_j} \sum_{i = 1}^s D_l[i,j] \norm{z_i} \\
         = & \begin{bmatrix}
            \norm{z_1} \\
            \vdots \\
            \norm{z_s}
            \end{bmatrix}^T D_l \begin{bmatrix}
            \norm{z_1} \\
            \vdots \\
            \norm{z_s}
            \end{bmatrix} \leq  
             \norm{\begin{bmatrix}
            \norm{z_1} \\
            \vdots \\
            \norm{z_s}
            \end{bmatrix}}^2 \lVert D_l \rVert_2 = \lVert x \rVert_2^2 \lVert D_l \rVert_2,  
    \end{align*}
    which shows the claim. 
\end{proof}

\end{document}
